\documentclass[10pt,a4paper]{amsart}
\usepackage[margin=19mm]{geometry}
\usepackage{amsmath,amssymb,mathtools}
\usepackage[scr=rsfs,cal=euler]{mathalfa}
\usepackage{xfrac}
\usepackage[unicode,hidelinks,bookmarks=false]{hyperref}
\newcommand{\ie}{i.e.\ }
\newcommand{\cf}{cf.\ }
\newcommand{\resp}{respectively\ }
\newcommand{\Subsection}{\subsection}
\providecommand{\nobreakdash}{\nobreak\hbox{-}}
\setlength{\emergencystretch}{2em}
\multlinegap0pt
\usepackage{amssymb}
\usepackage[normalem]{ulem}
\usepackage{tikz}
\usepackage[shortlabels]{enumitem}
\newtheorem{lemma}{Lemma}
\newtheorem{prop}{Proposition}
\title{The elementary theory of free Steiner triple systems}
\author{Silvia Barbina, Enrique Casanovas}
\date{Source: arXiv:2411.13723v3 (CC BY 4.0)}
\begin{document}
\maketitle

\noindent\emph{Source and licence.} The elementary theory of free Steiner triple systems. Version arXiv:2411.13723v3 (CC BY 4.0), distributed by arXiv under the Creative Commons Attribution 4.0 International licence, http://creativecommons.org/licenses/by/4.0/, as recorded in the arXiv licence metadata for this exact version and shown on the abstract page https://arxiv.org/abs/2411.13723v3. Full licence notice: \texttt{licenses/arxiv-2411.13723-CC-BY-4.0.txt}.\par\smallskip

\noindent\textit{Editorial scope.} Counted unit: the article's prop tree3 (source0.tex lines 841--846). The article's complete original proof of it is delivered with it (source0.tex lines 847--859, one \texttt{proof} environment of 13 lines). No mathematics is written, repaired or re-derived: the statement and the proof are the article's own lines, in the article's own notation, and no retained line is substituted or re-flowed.  Retained uncounted as the source-local context the proof needs: lines 829--834.  Nothing was omitted from any retained span, so the package carries no omission record.  No retained line contains a citation, so the wrapper carries no bibliography and no bibliography entry.  No retained line contains a figure environment, an image-inclusion command or drawing code, so no image is delivered and the package declares no asset dependency.  Editorial formatting only, in addition: an AMS \texttt{amsart} preamble that restates the article's own declarations and macros used by the retained lines, namely the environments \texttt{lemma}, \texttt{prop} on separate counters, together with the standard packages those lines need.  The retained environments print in this document as Lemma 1, Proposition 1 in order of appearance, whereas the article prints the same items with the numbering of its own full document.  The complete native source of the article as distributed in the arXiv e-print for this exact version is bundled unchanged as \texttt{sources/168843/source0.tex}.
\begin{lemma}\label{tree2}
	\begin{enumerate}
		\item An HF-ordering of an infinite STS does not have a maximum.
		\item If $<$ is an HF-ordering of an infinite Steiner triple system $M$, then for every $a\in M$, for every  $b \geq a$, there are $b_1,b_2>b$ such that  $a=b_1\cdot b_2$.
	\end{enumerate}
\end{lemma}

\begin{prop}\label{tree3} Let $<$ be an HF-ordering of the infinite Steiner triple system $M$ and let $n<\omega$. For every $a\in M$    there is a binary tree $T=\{a_s\mid s\in 2^{<n}\}$  in $M$  such that
	\begin{enumerate}
		\item $a_s>a$  for every $s$
		\item $a_s>a_t$  if  $s\subsetneq t$.
	\end{enumerate}
\end{prop}

\begin{proof}
%	By induction on $n$. The initial cases $n=0,1$ are trivial. Consider the case  $n=2$. Let $b>a$. We use Lemma~\ref{tree2} to get some $b_1 >b_2 >b$ in $M$ such that $b=b_1\cdot b_2$. We put $a_\emptyset = b_1$,  $a_0=b_2$ and $a_1 = b$. 
 

%	Now we assume inductively the result holds for $n\geq 2$.
%	Given $a\in M$, the inductive hypothesis gives some tree $T_0=\{a_s\mid s\in 2^{<n}\}$ as required for $a$.  Again by the inductive hypothesis we can find in $M$ a tree $T_1=\{a^\prime_s\mid s\in 2^{<n}\}$ with $a^\prime_s>a_\emptyset$ and condition 2 holds. Now we combine $T_0$ and $T_1$ to get a binary tree $T=\{b_s \mid s\in 2^{<n+1}\}$. We set $b_{0s}= a_s$, $b_{1s}= a^\prime_s$  and  $b_\emptyset = a_\emptyset\cdot a^\prime_\emptyset$.  It is enough to check that   $a_\emptyset\cdot a^\prime_\emptyset > a^\prime_\emptyset$. Since $n\geq 2$ this is clear, since otherwise the blocks $\{a^\prime_\emptyset,a^\prime_0,a^\prime_1\}$  and  $\{a^\prime_\emptyset, a_\emptyset, a_\emptyset \cdot a^\prime_\emptyset\}$ are different and contradict that $<$ is an HF-ordering.
	We prove by induction on $n$ that  for every $a\in M$ for every finite $A\subseteq M$, there is a binary tree $T=\{a_s\mid s\in 2^{<n}\}$  in $M$  satisfying conditions 1,2 and such that there is no block connecting elements of $T$ with elements of $A$ (that is, no products $t \cdot a$ with $t \in T$ and $a \in A$ are in $T \cup A$).
	 The initial cases $n=0,1$ are trivial. Consider the case  $n=2$. Let $A^\prime$ be the set of products $\{x\cdot y\mid x,y\in A\}$, and let $b>a$ and larger than every element of $A^\prime$. We use Lemma~\ref{tree2} to get some $b_1 >b_2 >b$ in $M$ such that $b=b_1\cdot b_2$. We let $a_\emptyset = b_1$,  $a_0=b_2$ and $a_1 = b$. 
	
	
	Now we assume inductively the result holds for $n\geq 2$.
	Given $a\in M$ and a finite set $A\subseteq M$, the inductive hypothesis gives some tree $T_0=\{a_s\mid s\in 2^{<n}\}$ as required for $a,A$.  Let $A^\prime =A\cup T_0$. Again by inductive hypothesis, we can find a tree $T_1=\{a^\prime_s\mid s\in 2^{<n}\}$ in $M$  with $a^\prime_s$ larger than every element of $A^\prime$, with no blocks connecting $T_1$ and $A^\prime$  and such that condition 2 holds. Now we combine $T_0$ and $T_1$ to get a binary tree $T=\{b_s \mid s\in 2^{<n+1}\}$. We set $b_{0s}= a_s$, $b_{1s}= a^\prime_s$  and  $b_\emptyset = a_\emptyset\cdot a^\prime_\emptyset$.  We first  check that   $a_\emptyset\cdot a^\prime_\emptyset > a^\prime_\emptyset$. Since $n\geq 2$ this is clear, since otherwise the blocks $\{a^\prime_\emptyset,a^\prime_0,a^\prime_1\}$  and  $\{a^\prime_\emptyset, a_\emptyset, a_\emptyset \cdot a^\prime_\emptyset\}$ are different and contradict that $<$ is an HF-ordering. Since there is no block connecting elements of $T_1$ with elements of $T_0$, $T$ is a tree.
\end{proof}

\end{document}
